% Generated draft: requires mathematical and source review.
\documentclass[11pt]{article}
\usepackage[margin=25mm]{geometry}
\usepackage{fontspec}
\setmainfont{Latin Modern Roman}
\usepackage{amsmath,amssymb,mathtools,unicode-math}
\setmathfont{Latin Modern Math}
\usepackage{enumitem,needspace,xurl,hyperref}
\setlength{\parindent}{0pt}
\setlength{\parskip}{5pt}
\setlength{\emergencystretch}{3em}
\newcommand{\N}{\mathbb{N}}
\newcommand{\Z}{\mathbb{Z}}
\newcommand{\Q}{\mathbb{Q}}
\newcommand{\R}{\mathbb{R}}
\newcommand{\C}{\mathbb{C}}
\newcommand{\F}{\mathbb{F}}
\newcommand{\A}{\mathbb{A}}
\newcommand{\PP}{\mathbb{P}}
\newcommand{\E}{\mathbb{E}}
\newcommand{\eps}{\varepsilon}
\newcommand{\dd}{\,\mathrm{d}}
\begin{document}
\section*{UnsolvedMath ID: 800607}
\textbf{Source problem code:} catalogue entry
\subsection{Definitions and mathematical statement}
The supplied record asks about three splitting properties but does not define its cover conventions or its ambient class of spaces. The following gives a parameterized formulation using conventional cover terminology; these missing conventions require source verification before this can be treated as a definitive reconstruction.

For a topological space \(Y\), a nontrivial cover is a family \(\mathcal U\subseteq\mathcal P(Y)\) such that \(\bigcup\mathcal U=Y\) and \(Y\notin\mathcal U\). Let \(\mathcal A(Y)\) denote the intended family of admissible nontrivial covers. The supplied record does not specify whether admissible covers must be open, clopen, or countable.

An \(\omega\)-cover is a nontrivial cover \(\mathcal U\) such that
\[
\forall F\subseteq Y\;\bigl(F\text{ finite}\Longrightarrow\exists U\in\mathcal U\; F\subseteq U\bigr).
\]
A cover is large if every point belongs to infinitely many of its members. A \(\tau\)-cover is a large nontrivial cover \(\mathcal U\) such that, for every \(x,y\in Y\), at least one of the sets
\[
\{U\in\mathcal U:x\in U,\ y\notin U\},\qquad
\{U\in\mathcal U:y\in U,\ x\notin U\}
\]
is finite. Write \(\Omega_{\mathcal A}(Y)\) and \(\mathrm T_{\mathcal A}(Y)\) for the admissible \(\omega\)-covers and admissible \(\tau\)-covers, respectively.

For two specified families of covers \(\mathcal B(Y)\) and \(\mathcal D(Y)\), use the conventional disjoint-subcover interpretation
\[
\mathsf{Split}(\mathcal B,\mathcal D)(Y)
\iff
\forall\mathcal U\in\mathcal B(Y)\;\exists\mathcal V,\mathcal W\in\mathcal D(Y)\;
\bigl(\mathcal V\subseteq\mathcal U\ \land\ \mathcal W\subseteq\mathcal U\ \land\ \mathcal V\cap\mathcal W=\varnothing\bigr).
\]
Here \(\mathcal V\) and \(\mathcal W\) are themselves covers of \(Y\); this interpretation does not require \(\mathcal V\cup\mathcal W=\mathcal U\). The intended definition of \(\mathsf{Split}\) must also be checked against the source.

For a fixed intended class \(\mathcal C\) of spaces and a fixed admissibility convention on these spaces and their powers, the question is whether any of
\[
P\in\left\{
\mathsf{Split}(\Omega_{\mathcal A},\Omega_{\mathcal A}),\quad
\mathsf{Split}(\Omega_{\mathcal A},\mathrm T_{\mathcal A}),\quad
\mathsf{Split}(\mathrm T_{\mathcal A},\mathrm T_{\mathcal A})
\right\}
\]
satisfies
\[
\forall X\in\mathcal C\;\left(P(X)\Longrightarrow
\forall n\in\{1,2,3,\ldots\}\;P(X^n)\right),
\]
where \(X^n\) has the product topology. Positive finite powers are used in this displayed formulation. The source record does not explicitly address the zeroth power. Neither \(\mathcal C\) nor \(\mathcal A\) is determined by the supplied excerpt, so the displayed question remains conditional on those choices.
\subsection{Short English statement}
Tsaban's Problem 35 asks whether any of the three properties \(\mathsf{Split}(\Omega,\Omega)\), \(\mathsf{Split}(\Omega,\mathrm T)\), and \(\mathsf{Split}(\mathrm T,\mathrm T)\) passes from a space to each of its positive finite Cartesian powers. The intended spaces, admissible covers, and exact splitting convention are missing from the supplied excerpt and require verification.
\subsection{Sources}
\begin{itemize}
\item Supplied local UnsolvedMath registry record: permanent ID \(800607\), exact source problem code \texttt{catalogue entry}; upstream ID \(800607\), revision \texttt{d371ee917ac9647fbb06cd6fc529deec6f3dda82}.
\item Supplied attribution: Tsaban, \emph{Selection Principles and Special Sets of Reals: Open Problems} (2006), Problem 35. Supplied source URL: \url{https://arxiv.org/abs/math/0609601}. The registry records access on 2026-10-16 and extraction by \texttt{arxiv\_latex}; these are supplied provenance details, not a claim of independent access or verification.
\item The supplied attribution/context marker is \texttt{[split]}; its bibliographic expansion is absent from the record.
\item The registry labels the problem \texttt{open} and gives an update date of 2026-10-16. This is a supplied status label, not an independently verified assessment of the mathematical literature. No separate literature-verification date or additional literature references are supplied.
\end{itemize}
\end{document}
